Kathmandu Valley's peak-hour arterial congestion has grown with motorization far
faster than road capacity: a 2012 survey found motorcycle traffic volumes had
increased roughly 19.8$\times$ since 1991 (a 16\% annual growth rate)
\cite{jica2012}, concentrated on a road network that has not grown proportionally.
Two interventions are available to a city without the years and capital a new road
or metro line requires: retiming existing traffic signals, and charging for road
use at the busiest times. Both are software- or policy-level changes to
infrastructure that already exists, in contrast to physical capacity expansion.

This paper asks two questions. First, do signal retiming and congestion pricing
each reduce delay significantly on Kathmandu's arterial network, and does
combining them interact super-additively -- a mechanism claim, not only an
effect-size claim, that would imply pricing relieves saturation enough for
adaptive signal control to regain room to act (Section~\ref{sec:core})? We find
the first holds and test the second directly, and report that it does not hold
statistically, rather than assume it from the combined effect's size alone.
Second, how do these results hold up against the honest
limits of a synthesized-demand simulation: against observed traffic counts
(Section~\ref{sec:validation}), against the model's own fitted-parameter
sensitivity (Section~\ref{sec:sensitivity}), and against four further candidate
interventions tested with the same rigor, two of which are rejected
(Section~\ref{sec:additional})?

We build a calibrated mesoscopic simulation of Kathmandu's Ring Road and
city-core arterial network (233 nodes, 590 links) and test every comparison across
15 paired random seeds, reporting a confidence interval rather than a single-run
point estimate -- a design choice motivated directly by observing that small
changes in this model's demand can move single-run delay by several percent on
noise alone (Section~\ref{sec:stats}). This same design is what let us catch that
an apparent super-additive interaction did not survive a properly computed
confidence interval (Section~\ref{sec:core}) -- a correction made during this
project's own internal review, reported here rather than silently dropped.
